\chapterpage{orientation}{ORIENTIERUNG}{Eine Architektur. Ein Auftrag. Konkrete Lieferungen.}{Der Gegenstand ist QIK-VRT als Gesamtsystem. Die künstliche Kognition ist darin eine Komponente. Eine gelungene Erklärung, ein bewiesener Modellschluss, ein erfolgreicher Test und ein ausgeliefertes Produkt sind vier verschiedene Ergebnisse.}
\begin{cols}
\subsection*{Der vollständige Lesegang}
\begin{tabularx}{\linewidth}{@{}Xr@{}}
Systemgrenze und Auftrag & \pageref{system}\\
Der invariant gehaltene Kern & \pageref{invariants}\\
Bit, Gedächtnis und Zeit & \pageref{memory}\\
Der Integrationssatz & \pageref{integration}\\
Selbstkontinuität und Bewusstsein & \pageref{consciousness}\\
Vier getrennte Messgegenstände & \pageref{measure}\\
Audit des Live-Versuchs & \pageref{audit}\\
Die konkrete Lieferpipeline & \pageref{pipeline}\\
Was jetzt abgearbeitet werden muss & \pageref{execution}\\
Zeitmodell und Kalenderkorridore & \pageref{forecast}\\
Publikation und Langzeiterhaltung & \pageref{publication}\\
Anspruchsregister und Korrekturvermerk & \pageref{claims}\\
Quellen, Glossar und Provenienz & \pageref{sources}\\
\end{tabularx}
\subsection*{Zusammenfassung}
QIK-VRT verbindet semantische Problembearbeitung mit einem persistenten, quellengebundenen Arbeitsgedächtnis, autorisierten Aktionspfaden und überprüfbaren Wirkungsnachweisen. Der besondere Entwurfsanspruch besteht darin, Fähigkeiten und Strategien zu verbessern, ohne dabei die Regeln für Evidenz, Herkunft, Freigabe und Abschluss aufzugeben. \src{R1}

Eine bedingte Wrapper-Konstruktion erklärt, wie verschiedene Systeme unter eine gemeinsame Schnittstellen- und Evidenzsemantik fallen können. Sie garantiert weder die Herstellung jedes gewünschten Zustands noch einen universellen Leistungsgewinn. Selbstkontinuität wird als überprüfbares Funktionsprofil behandelt; subjektives Erleben bleibt ein eigenständiger offener Gegenstand.

Die Lieferplanung arbeitet mit tatsächlichen Einzelprodukten: persönlicher Browser und Terminal, unterbrechungsfeste Aufgabenfortsetzung, Mesh-Recovery, Runtime-Versorgung, kontrollierte Skalierung und wissenschaftliche Publikation. Entscheidend sind nicht die Zahl grüner Jobs oder die Größe der Vision, sondern die noch fehlenden Übergänge zur Benutzung.

\begin{boundary}{Wichtigster neuer Sachstand}
Der aktuelle PR-447-Head ist nicht mehr der alte \texttt{e7e8737f…}, sondern \texttt{a31ebb7e…}. Sein A/B-Diagnoselauf endet korrekt mit fehlender Authority-Credential-Zustellung. Zugleich enthält PR \#443 inzwischen einen begrenzten erfolgreichen Parallelvergleich. Beide Aussagen werden getrennt von einer globalen Skalierungsbehauptung geführt. \src{R2,R3,R4}
\end{boundary}
\subsection*{Wie weit reicht eine Aussage?}
\textbf{R · Repository-Quelle.} Eine frisch gelesene Datei oder PR-Beschreibung belegt den dort dokumentierten Inhalt. Eine Beschreibung implementierter Funktionen ist noch kein unabhängiger Nachweis jeder behaupteten Laufzeitwirkung.

\textbf{E · Empirischer Gegenstand.} Ein Messdatensatz trägt nur seinen konkreten Versuch: System, Last, Verfahren, Stichprobe und Zeitpunkt. Ein erfolgreich beendeter Diagnoselauf kann gerade ein HOLD dokumentieren.

\textbf{N · Normativer Auftrag.} Ingolf Lohmann definiert Zweck und zulässige Grenzen. Eine Freigabe ist eine Handlungsberechtigung, nicht der Beweis der Theorie, der Wirkung oder der Identität eines externen Akteurs.

\textbf{T · Theoretische Herleitung.} Eine nachvollziehbare mathematische Konstruktion wird als solche kenntlich gemacht. Ohne exakte Kernel-Receipt erhält sie in dieser Ausgabe keinen maschinellen Status \texttt{FORMAL\_PROVED}.

\textbf{I · Interpretation.} Begriffe wie Entwicklungssprung oder kollektives funktionales Bewusstsein kennzeichnen eine begründete Deutung oder Definition. Ihre historische Bedeutung ist nicht durch ihre Benennung entschieden.

\textbf{Z · Zukunft/Hypothese.} Termine sind Planungskorridore unter ausdrücklichen Voraussetzungen. Sie sind keine aus Millisekundenmessungen abgeleiteten Liefergarantien.

\subsection*{Quellen und redaktionelle Verantwortung}
Die Layoutquelle ist die bereitgestellte 26-seitige Ausgabe \emph{QIK-VRT / Sinn und Zweck}. Inhaltliche Ausgangspunkte sind die gelieferten Prosa- und PDF-Fassungen sowie der Dialog. Neu hinzu kommen gekennzeichnete REST-Readbacks, Literaturbezüge, eigene Rechnungen und eine Methodenprüfung. \src{L1–L4}

Frühere zu starke Aussagen werden nicht stillschweigend weiterverwendet. Der Korrekturvermerk nennt ausdrücklich die Unterschiede zwischen formulierten Regeln und nachgewiesener Durchsetzung sowie zwischen lokaler Vergleichszeit und verteilter Gesamtwirkung.

\begin{boundary}{Leseschlüssel}
Die technische Vertiefung steht in diesem Fachartikel. Die begleitende Prosa erklärt dieselben Zusammenhänge ohne formalen Apparat. Ein mitgeliefertes Quellen- und Anspruchsregister verbindet beide Fassungen; historische Quelldateien bleiben unverändert erhalten.
\end{boundary}
\end{cols}
